\chapter{Álgebra de límites y propiedades de orden}

\section{Producto de sucesiones}

Habiendo establecido el teorema del límite algebraico para sumas y múltiplos
escalares, ahora abordamos el producto de sucesiones convergentes. Este resultado,
combinado con el trabajo anterior, completa la estructura algebraica de los límites:
podemos sumar, multiplicar por constantes, multiplicar sucesiones entre sí e incluso
dividir (siempre que el límite sea no nulo).

La estrategia de demostración refleja la regla del producto del cálculo, donde
$(fg)' = f'g + g'f$. Introducimos artificialmente términos que nos permiten separar
las contribuciones de cada sucesión.

\subsection{El enunciado matemático}

\begin{theorem}[ProdLimNeNe]
  Si $a, b : \mathbb{N} \to \mathbb{R}$ son sucesiones con $a(n) \to L$ y
  $b(n) \to M$, donde $L \neq 0$ y $M \neq 0$, y $c(n) = a(n) \cdot b(n)$
  para todo $n$, entonces:
  \[
    c(n) \to L \cdot M
  \]
\end{theorem}

\subsection{Enfoque estratégico}

La idea clave es escribir:
\[
  a(n) \cdot b(n) - L \cdot M
  = \bigl(a(n) - L\bigr) \cdot b(n) + L \cdot \bigl(b(n) - M\bigr)
\]

Esta descomposición separa el error en dos partes:
\begin{itemize}
  \item $\bigl(a(n)-L\bigr)\cdot b(n)$: el error en $a$ multiplicado por el valor de $b$.
  \item $L\cdot\bigl(b(n)-M\bigr)$: el límite $L$ multiplicado por el error en $b$.
\end{itemize}

Para acotar el primer término usamos que $b$ es convergente y por tanto acotada
(por algún $K>0$). Entonces $|a(n)-L|$ puede hacerse menor que $\varepsilon/(2K)$,
lo que nos da control sobre ese término.

Para el segundo término usamos directamente que $b(n)\to M$, haciendo
$|b(n)-M|$ menor que $\varepsilon/(2|L|)$.

\subsection{Demostración en Lean}

\begin{minted}{lean}
-- Demostración en lenguaje natural
-- ================================

-- Sea ε > 0. Como aₙ converge a L, existe N₁ tal que si n ≥ N₁,
-- entonces
--    |aₙ − L| < 1
-- Por tanto, para n ≥ N₁
--    |aₙ| = |aₙ − L + L|
--         ≤ |aₙ − L| + |L|
--         < 1 + |L|
-- Definamos
--    C = 1 + |L|
-- Entonces
--    C > 0
-- y, para n ≥ N₁
--    |aₙ| ≤ C
-- Ahora, como aₙ converge a L, existe N₂ tal que si n ≥ N₂,
-- entonces
--    |aₙ − L| < ε/(2(|M| + 1))
-- Por otra parte, como bₙ converge a M, existe N₃ tal que si
-- n ≥ N₃, entonces
--    |bₙ − M| < ε/(2C)
-- Sea
--    N = máx⁡{N₁,N₂,N₃}
-- Si n ≥ N, entonces se cumplen las tres desigualdades
-- anteriores. Ahora estimamos:
--    |aₙbₙ − LM|
-- Sumamos y restamos aₙM:
--    aₙbₙ − LM = aₙbₙ − aₙM + aₙM − LM.
-- Entonces
--    |aₙbₙ − LM| = |aₙ(bₙ − M) + M(aₙ − L)|
-- Por la desigualdad triangular,
--    |aₙbₙ − LM| ≤ |aₙ||bₙ − M| +|M||aₙ − L|
-- Como n ≥ N₁,tenemos
--    |aₙ| ≤ C
-- Luego,
--    |aₙbₙ − LM| ≤ C|bₙ − M| + |M||aₙ − L|
-- Usando las cotas elegidas,
--    C|bₙ − M| < C⋅ε/(2C)
--              = ε/2
-- y
--    |M||aₙ − L| < |M|⋅ε/(2(|M| + 1))
-- Como
--    |M|/(|M| + 1) < 1
-- se sigue que
--    |M|⋅ε/(2(|M| + 1)) < ε/2
-- Por tanto,
--    |aₙbₙ − LM| < ε/2 + ε/2
--                = ε
-- Así, para todo ε > 0 existe N tal que si n ≥ N, entonces
--    |aₙbₙ − LM| < ε
-- Esto prueba que, para todo ε > 0, existe N tal que |aₙbₙ − LM| < ε
-- siempre que n ≥ N, es decir, aₙbₙ converge a LM.

-- Demostraciones en Lean 4
-- ========================

import Mathlib.Data.Real.Basic
import Analisis_en_Lean.C2_El_calculo_de_pi
import Analisis_en_Lean.C9B_Acotacion_de_convergentes
import Mathlib.Tactic

variable {a b c : ℕ → ℝ}
variable {L M : ℝ}

-- 1ª solución
-- ===========

example
  (ha : LimSuc a L)
  (hb : LimSuc b M)
  (hc : ∀ n, c n = a n * b n)
  : LimSuc c (L * M) :=
by
  intro ε hε
  -- ε : ℝ
  -- hε : ε > 0
  -- ⊢ ∃ k, ∀ n ≥ k, |c n - L * M| < ε
  obtain ⟨N₁, hN₁⟩ := ha 1 (by positivity)
  -- N₁ : ℕ
  -- hN₁ : ∀ n ≥ N₁, |a n - L| < 1
  have h1 : ∀ n ≥ N₁, |a n| < 1 + |L| := by grind
  set C := 1 + |L|
  -- h1 : ∀ n ≥ N₁, |a n| < C
  have h2 : C > 0 := by grind
  obtain ⟨N₂, hN₂⟩ := ha (ε / (2 * (|M| + 1))) (by positivity)
  -- N₂ : ℕ
  -- hN₂ : ∀ n ≥ N₂, |a n - L| < ε / (2 * (|M| + 1))
  obtain ⟨N₃, hN₃⟩ := hb (ε / (2 * C)) (by positivity)
  -- N₃ : ℕ
  -- hN₃ : ∀ n ≥ N₃, |b n - M| < ε / (2 * C)
  set N := max N₁ (max N₂ N₃)
  use N
  -- ⊢ ∀ n ≥ N, |c n - L * M| < ε
  intro n hn
  -- n : ℕ
  -- hn : n ≥ N
  -- ⊢ |c n - L * M| < ε
  have h3 : n ≥ N₁ := by grind
  have h4 : n ≥ N₂ := by grind
  have h5 : n ≥ N₃ := by grind
  have h6 : |a n| < C := by grind
  have h7 : |b n - M| < ε / (2 * C) := by grind
  have h8 : |a n - L| < ε / (2 * (|M| + 1)) := by grind
  have h9 : |M| * (ε / (2 * (|M| + 1))) < ε / 2 := by
              field_simp
              -- ⊢ |M| < |M| + 1
              norm_num
  calc |c n - L * M|
       = |a n * b n - L * M|                             := by grind
     _ = |a n * (b n - M) + M * (a n - L)|               := by grind
     _ ≤ |a n * (b n - M)| + |M * (a n - L)|             := by grind
     _ = |a n| * |b n - M| + |M| * |a n - L|             := by grind
     _ ≤ C * |b n - M| + |M| * |a n - L|                 := by gcongr
     _ < C * (ε / (2 * C)) + |M| * |a n - L|             := by gcongr
     _ ≤ C * (ε / (2 * C)) + |M| * (ε / (2 * (|M| + 1))) := by gcongr
     _ = ε / 2 + |M| * (ε / (2 * (|M| + 1)))             := by grind
     _ < ε / 2 + ε / 2                                   := by grind
     _ = ε                                               := by grind

-- 2ª solución
-- ===========

lemma L1
  (h : LimSuc a L)
  : ∃ k, ∀ n ≥ k, |a n| < 1 + |L| :=
by
  obtain ⟨k, hk⟩ := h 1 one_pos
  -- k : ℕ
  -- hk : ∀ n ≥ k, |a n - L| < 1
  use k
  -- ⊢ ∀ n ≥ k, |a n| < |L| + 1
  intro n hn
  -- n : ℕ
  -- hn : n ≥ k
  -- ⊢ |a n| < |L| + 1
  calc |a n|
       = |(a n - L) + L| := congrArg abs (sub_add_cancel (a n) L).symm
     _ ≤ |a n - L| + |L| := abs_add_le (a n - L) L
     _ < 1 + |L|         := (add_lt_add_iff_right |L|).mpr (hk n hn)

example
  (ha : LimSuc a L)
  (hb : LimSuc b M)
  (hc : ∀ n, c n = a n * b n)
  : LimSuc c (L * M) :=
by
  intro ε hε
  -- ε : ℝ
  -- hε : ε > 0
  -- ⊢ ∃ k, ∀ n ≥ k, |c n - L * M| < ε
  obtain ⟨N₁, hN₁⟩ := L1 ha
  -- N₁ : ℕ
  -- hN₁ : ∀ n ≥ N₁, |a n| < 1 + |L|
  set C := 1 + |L| with hC
  -- hC : C = 1 + |L|
  have h2 : C > 0 := by positivity
  obtain ⟨N₂, hN₂⟩ := ha (ε / (2 * (|M| + 1))) (by positivity)
  -- N₂ : ℕ
  -- hN₂ : ∀ n ≥ N₂, |a n - M| < ε / (2 * (|M| + 1))
  obtain ⟨N₃, hN₃⟩ := hb (ε / (2 * C)) (by positivity)
  -- N₃ : ℕ
  -- hN₃ : ∀ n ≥ N₃, |b n - M| < ε / (2 * C)
  set N := max N₁ (max N₂ N₃)
  use N
  -- ⊢ ∀ n ≥ N, |c n - L * M| < ε
  intro n hn
  -- n : ℕ
  -- hn : n ≥ N
  -- ⊢ |c n - L * M| < ε
  have h3 : n ≥ N₁ := by omega
  have h4 : n ≥ N₂ := by omega
  have h5 : n ≥ N₃ := by omega
  have h6 : |a n| < C := by
    calc |a n|
         < 1 + |L| := hN₁ n h3
       _ = C       := rfl
  have h7 : |b n - M| < ε / (2 * C) := hN₃ n h5
  have h8 : |a n - L| < ε / (2 * (|M| + 1)) := hN₂ n h4
  have h9 : |M| * (ε / (2 * (|M| + 1))) < ε / 2 := by
              field_simp
              -- ⊢ |M| < |M| + 1
              norm_num
  calc |c n - L * M|
       = |a n * b n - L * M| :=
            congrArg (|· - L * M|) (hc n)
     _ = |a n * (b n - M) + M * (a n - L)| :=
            by ring_nf
     _ ≤ |a n * (b n - M)| + |M * (a n - L)| :=
            abs_add_le (a n * (b n - M)) (M * (a n - L))
     _ = |a n| * |b n - M| + |M| * |a n - L| :=
            by simp only [abs_mul]
     _ ≤ C * |b n - M| + |M| * |a n - L| :=
            by gcongr
     _ < C * (ε / (2 * C)) + |M| * |a n - L| :=
            by gcongr
     _ ≤ C * (ε / (2 * C)) + |M| * (ε / (2 * (|M| + 1))) :=
            by gcongr
     _ = ε / 2 + |M| * (ε / (2 * (|M| + 1))) :=
            by field_simp
     _ < ε / 2 + ε / 2 :=
            by gcongr
     _ = ε :=
            add_halves ε
\end{minted}

\subsection{Demostración en lenguaje natural}

  Sea $\varepsilon > 0$ dado. Como $b(n) \to M$ con $M \neq 0$, por
  \texttt{BddOfConvNonzero} existe $K > 0$ tal que $|b(n)| \leq K$ para todo $n$.

  Definamos:
  \[
    \varepsilon_1 := \frac{\varepsilon}{2K}, \qquad
    \varepsilon_2 := \frac{\varepsilon}{2|L|}
  \]

  Nótese que $\varepsilon_1 > 0$ (pues $K > 0$) y $\varepsilon_2 > 0$
  (pues $L \neq 0$ implica $|L| > 0$).

  Como $a(n) \to L$, existe $N_1$ tal que para todo $n \geq N_1$:
  \[
    |a(n) - L| < \varepsilon_1 = \frac{\varepsilon}{2K}
  \]

  Como $b(n) \to M$, existe $N_2$ tal que para todo $n \geq N_2$:
  \[
    |b(n) - M| < \varepsilon_2 = \frac{\varepsilon}{2|L|}
  \]

  Sea $N := N_1 + N_2$. Para cualquier $n \geq N$ tenemos $n \geq N_1$ y
  $n \geq N_2$.

  Usando la identidad algebraica
  $a(n)\cdot b(n) - L\cdot M = (a(n)-L)\cdot b(n) + L\cdot(b(n)-M)$
  y aplicando la desigualdad triangular:
  \begin{align*}
    |c(n) - L\cdot M|
    &= |(a(n)-L)\cdot b(n) + L\cdot(b(n)-M)| \\
    &\leq |(a(n)-L)\cdot b(n)| + |L\cdot(b(n)-M)| \\
    &= |a(n)-L|\cdot|b(n)| + |L|\cdot|b(n)-M|
  \end{align*}

  Para el primer término, como $|b(n)| \leq K$ y $|a(n)-L| < \varepsilon/(2K)$:
  \[
    |a(n)-L|\cdot|b(n)| < \frac{\varepsilon}{2K}\cdot K = \frac{\varepsilon}{2}
  \]

  Para el segundo término, como $|b(n)-M| < \varepsilon/(2|L|)$:
  \[
    |L|\cdot|b(n)-M| < |L|\cdot\frac{\varepsilon}{2|L|} = \frac{\varepsilon}{2}
  \]

  Por tanto:
  \[
    |c(n) - L\cdot M| < \frac{\varepsilon}{2} + \frac{\varepsilon}{2} = \varepsilon
  \]
  \hfill$\square$

\subsection{Conexión con el cálculo: ¿por qué sumar y restar \texorpdfstring{$b(n)\cdot L$}{b(n)·L}?}

La técnica de la demostración refleja la regla del producto del cálculo:
$(fg)' = f'g + g'f$. Para entender por qué sumamos y restamos $b(n)\cdot L$,
pensemos geométricamente en el área de un rectángulo (al fin y al cabo, la
multiplicación es rectángulos).

Cuando $a(n)$ y $b(n)$ están cerca de sus límites $L$ y $M$, el producto
$a(n)\cdot b(n)$ representa el área de un rectángulo de lados $a(n)$ y $b(n)$,
mientras que $L\cdot M$ es el área del rectángulo límite de lados $L$ y $M$.

La diferencia de áreas satisface:
\[
  a(n)\cdot b(n) - L\cdot M
  = L\cdot(b(n)-M) + M\cdot(a(n)-L) + (a(n)-L)\cdot(b(n)-M)
\]

La clave es reemplazar $M$ por $b(n)$ en el término central (sumando y restando
$b(n)\cdot L$) para obtener:
\[
  a(n)\cdot b(n) - L\cdot M
  = (a(n)-L)\cdot b(n) + L\cdot(b(n)-M)
\]

¿Por qué es mejor esto? Porque ahora cada término aísla una fuente de error:
\begin{itemize}
  \item $(a(n)-L)\cdot b(n)$: error pequeño en $a$ por la cantidad acotada $b(n)$.
  \item $L\cdot(b(n)-M)$: la constante fija $L$ por el error pequeño en $b$.
\end{itemize}

\subsection{Completando el teorema del límite algebraico}

Este resultado completa el teorema del límite algebraico. Combinado con los
resultados anteriores sobre sumas, múltiplos escalares y recíprocos de sucesiones
con límite no nulo, podemos calcular el límite de cualquier combinación algebraica
de sucesiones convergentes.

Por ejemplo, si $a(n)\to L$ y $b(n)\to M$, el límite de
\[
  \frac{a(n)^2 + 2a(n) + b(n)}{3b(n) + 2 - a(n)^2}
\]
es, siempre que el límite del denominador sea no nulo:
\[
  \frac{L^2 + 2L + M}{3M + 2 - L^2}
\]

\section{Paso al límite en desigualdades}

El paso al límite en desigualdades establece una relación fundamental entre
desigualdades y límites: si una sucesión está acotada superiormente (o
inferiormente) por una constante, entonces su límite también lo está. Este
resultado respeta la estructura de orden de los números reales.

\subsection{Nuevas herramientas}

\begin{definicion}[SeqBddBy]
  Una sucesión $a : \mathbb{N} \to \mathbb{R}$ está \emph{acotada por $M$}
  (denotado $\mathrm{SeqBddBy}(a, M)$) si $a(n) \leq M$ para todo
  $n \in \mathbb{N}$:
  \[
    \mathrm{SeqBddBy}(a,M) \iff \forall n\in\mathbb{N},\; a(n)\leq M
  \]

  Nótese la diferencia entre $\mathrm{SeqBdd}$ (acotada: $|a(n)|\leq M$) y
  $\mathrm{SeqBddBy}$ (acotada superiormente: $a(n)\leq M$).
\end{definicion}

\subsection{El enunciado matemático}

\begin{theorem}[OrderLimLe]
  Si $a:\mathbb{N}\to\mathbb{R}$ converge a $L$ y $a(n)\leq K$ para todo $n$,
  entonces $L\leq K$.
\end{theorem}

\subsection{Enfoque estratégico}

Lo demostramos por contradicción. Supongamos $L > K$. Entonces $L-K>0$, y
podemos usar esta cantidad positiva como $\varepsilon$ en la definición de
convergencia.

Por convergencia, existe $N$ tal que para todo $n\geq N$:
\[
  |a(n)-L| < L-K
\]
Esto implica $a(n) > K$, lo que contradice la hipótesis $a(n)\leq K$.

\subsection{Demostración en Lean}

\begin{minted}{lean}
-- Demostración en lenguaje natural
-- ================================

-- Lo demostramos por contradicción. Supongamos que
--    L > M                                                    (1)
-- y llegaremos a la contradicción M < M.
--
-- De (1), se tiene
--    L - M > 0
-- y, usando esta cantidad en la convergencia de aₙ, se obtiene
-- un k tal que
--    ∀ n ≥ k, |aₙ - L| < L - M                                (2)
-- En particular, para n = k (que satisface k ≥ k), de (2) se
-- obtiene
--    |aₖ - L| < L - M                                         (3)
-- Además,
--    L - aₖ ≤ |aₖ - L|                                        (4)
-- ya que
--    L - aₖ = -(aₖ - L)
--           ≤ |aₖ - L|
-- Finalmente,
--    M = L - (L - M)
--      < L - |aₖ - L|  [por (3)]
--      ≤ L - (L - aₖ)  [por (4)]
--      = aₖ
--      ≤ M             [porque M es cota superior de aₙ]
-- Por tanto, M < M.

-- Demostraciones en Lean 4
-- ========================

import Mathlib.Data.Real.Basic
import Analisis_en_Lean.C2_El_calculo_de_pi
import Mathlib.Tactic

def CotaSup (a : ℕ → ℝ) (M : ℝ) : Prop :=
  ∀ n, a n ≤ M

variable {a : ℕ → ℝ}
variable {L M : ℝ}

-- 1ª solución
-- ===========

example
  (ha : LimSuc a L)
  (hM : CotaSup a M)
  : L ≤ M :=
by
  by_contra hL
  -- hL : ¬L ≤ M
  -- ⊢ False
  apply lt_irrefl M
  -- ⊢ M < M
  obtain ⟨k, hk⟩ := ha (L - M) (by grind)
  -- k : ℕ
  -- hk : ∀ n ≥ k, |a n - L| < L - M
  calc M
       < a k := by grind
     _ ≤ M   := hM k

-- 2ª solución
-- ===========

example
  (ha : LimSuc a L)
  (hM : CotaSup a M)
  : L ≤ M :=
by
  by_contra hL
  -- hL : ¬L ≤ M
  -- ⊢ False
  apply lt_irrefl M
  -- ⊢ M < M
  have h1 : L > M := not_le.mp hL
  have h2 : L - M > 0 := sub_pos.mpr h1
  obtain ⟨k, hk⟩ := ha (L - M) h2
  -- k : ℕ
  -- hk : ∀ n ≥ k, |a n - L| < L - M
  calc M
       = L - (L - M)    := by grind
     _ < L - |a k - L|  := by grind
     _ ≤ L - (L - a k)  := by grind
     _ = a k            := by grind
     _ ≤ M              := hM k

-- 3ª solución
-- ===========

example
  (ha : LimSuc a L)
  (hM : CotaSup a M)
  : L ≤ M :=
by
  by_contra hL
  -- hL : ¬L ≤ M
  -- ⊢ False
  apply lt_irrefl M
  -- ⊢ M < M
  have h1 : L - M > 0 := sub_pos.mpr (not_le.mp hL)
  obtain ⟨k, hk⟩ := ha (L - M) h1
  -- k : ℕ
  -- hk : ∀ n ≥ k, |a n - L| < L - M
  have h2 : |a k - L| < L - M := hk k (le_refl k)
  have h3 : L - a k ≤ |a k - L| := by
    calc L - a k
         = -(a k - L) := by ring
       _ ≤ |a k - L|  := neg_le_abs (a k - L)
  calc M
       = L - (L - M)    := by ring
     _ < L - |a k - L|  := sub_lt_sub_left h2 L
     _ ≤ L - (L - a k)  := by gcongr
     _ = a k            := by ring
     _ ≤ M              := hM k

-- 4ª solución
-- ===========

example
  (ha : LimSuc a L)
  (hM : CotaSup a M)
  : L ≤ M :=
by
  by_contra hL
  -- hL : ¬L ≤ M
  -- ⊢ False
  apply lt_irrefl M
  -- ⊢ M < M
  obtain ⟨k, hk⟩ := ha (L - M) (sub_pos.mpr (not_le.mp hL))
  -- k : ℕ
  -- hk : ∀ n ≥ k, |a n - L| < L - M
  have h1 : |a k - L| < L - M := hk k (le_refl k)
  have h2 : L - a k ≤ |a k - L| := by
    calc L - a k
         = -(a k - L) := (neg_sub (a k) L).symm
       _ ≤ |a k - L|  := neg_le_abs (a k - L)
  calc M
       = L - (L - M)    := (sub_sub_self L M).symm
     _ < L - |a k - L|  := sub_lt_sub_left h1 L
     _ ≤ L - (L - a k)  := sub_le_sub_left h2 L
     _ = a k            := sub_sub_self L (a k)
     _ ≤ M              := hM k

-- Solución del libro
-- ==================

example
  (ha : LimSuc a L)
  (hM : CotaSup a M)
  : L ≤ M :=
by
  by_contra hL
  -- hL : ¬L ≤ M
  -- ⊢ False
  have hL' : M < L := by bound
  have hLM : 0 < (L - M) := by linarith [hL']
  obtain ⟨k, hk⟩ := ha (L - M) hLM
  -- k : ℕ
  -- hk : ∀ n ≥ k, |a n - L| < L - M
  specialize hk k (by bound)
  -- hk : |a k - L| < L - M
  rewrite [abs_lt] at hk
  -- hk : -(L - M) < a k - L ∧ a k - L < L - M
  have f1 : L - (L - M) < a k := by linarith [hk.1]
  have f2 : M = L - (L - M) := by ring_nf
  specialize hM k
  -- hM : a k ≤ M
  linarith [f2, hM, f1]

-- Lemas usados
-- ============

variable (x y : ℝ)
#check (le_refl x : x ≤ x)
#check (lt_irrefl x : ¬x < x)
#check (neg_le_abs x : -x ≤ |x|)
#check (neg_sub x y : -(x - y) = y - x)
#check (not_le : ¬x ≤ y ↔ y < x)
#check (sub_le_sub_left : x ≤ y → ∀ z, z - y ≤ z - x)
#check (sub_lt_sub_left : x < y → ∀ z, z - y < z - x)
#check (sub_pos : 0 < x - y ↔ y < x)
#check (sub_sub_self x y : x - (x - y) = y)
\end{minted}

\subsection{Demostración en lenguaje natural}

  Procedemos por contradicción. Supongamos $L > K$.

  Entonces $L-K > 0$. Como $a(n)\to L$, por la definición de convergencia con
  $\varepsilon = L-K > 0$, existe $N$ tal que para todo $n\geq N$:
  \[
    |a(n)-L| < L-K
  \]

  En particular, tomando $n=N$:
  \[
    |a(N)-L| < L-K \implies -(L-K) < a(N)-L < L-K
  \]

  De la desigualdad izquierda:
  \[
    a(N) > L - (L-K) = K
  \]

  Pero por hipótesis $a(N)\leq K$, lo que es una contradicción.

  Por tanto, $L > K$ es falso, y concluimos $L\leq K$. \hfill$\square$

\subsection{Importancia de la preservación del orden}

Este teorema establece que los límites preservan las desigualdades débiles.
Nótese que las desigualdades estrictas no se preservan necesariamente: si
$a(n)<K$ para todo $n$, solo podemos concluir $L\leq K$, no $L<K$.

Por ejemplo, la sucesión $a(n)=1/n$ satisface $a(n)>0$ para todo $n$, pero
$\lim_{n\to\infty}a(n)=0$, que no es estrictamente positivo.

El paso al límite en desigualdades se usa para demostrar:
\begin{itemize}
  \item El teorema del emparedado.
  \item El teorema de convergencia monótona.
  \item Diversos criterios de comparación para series.
  \item Propiedades del supremo e ínfimo de sucesiones.
\end{itemize}

\section{Subsucesiones}

Una subsucesión se forma seleccionando términos de una sucesión preservando
su orden. Hay dos perspectivas equivalentes. La mayoría de los autores llaman
$a_{\sigma(n)}$ «subsucesión de $a$», viéndola como una nueva sucesión extraída
de la original. Sin embargo, preferimos llamar $\sigma$ la subsucesión: es la
función estrictamente creciente $\sigma:\mathbb{N}\to\mathbb{N}$ que indica qué
términos conservar. La subsucesión de $a$ determinada por $\sigma$ es entonces
la composición $a\circ\sigma$.

Geométricamente, imagine la gráfica de la sucesión original $a(n)$. Para
formar una subsucesión se eliminan algunos términos y se reindexan los
restantes. Si conservamos solo los términos en las posiciones $0, 2, 5, 7, 11,\ldots$,
los reetiquetamos como $0, 1, 2, 3, 4,\ldots$ La función $\sigma$ codifica
este reetiquetado: $\sigma(0)=0$, $\sigma(1)=2$, $\sigma(2)=5$, etc.

El resultado fundamental es que si una sucesión converge, toda subsucesión
converge al mismo límite.

\subsection{Nuevas herramientas}

\begin{definicion}[Subseq]
  Una función $\sigma:\mathbb{N}\to\mathbb{N}$ es una \emph{subsucesión}
  (denotado $\mathrm{Subseq}(\sigma)$) si es estrictamente creciente:
  \[
    \mathrm{Subseq}(\sigma) \iff \forall\, i,j\in\mathbb{N},\;
    i < j \implies \sigma(i) < \sigma(j)
  \]

  La subsucesión de $a$ determinada por $\sigma$ es la composición
  $a\circ\sigma:\mathbb{N}\to\mathbb{R}$, donde $(a\circ\sigma)(n)=a(\sigma(n))$.
\end{definicion}

\begin{lemma}[SubseqGe]
  Si $\sigma$ es una subsucesión, entonces $n\leq\sigma(n)$ para todo $n$.
  Esto puede demostrarse por inducción y captura la intuición de que estamos
  «expandiendo» los índices.
\end{lemma}

\subsection{El enunciado matemático}

\begin{theorem}[SubseqConv]
  Si $a:\mathbb{N}\to\mathbb{R}$ converge a $L$ y $\sigma:\mathbb{N}\to\mathbb{N}$
  es una subsucesión, entonces $a\circ\sigma$ también converge a $L$.
\end{theorem}

\subsection{Enfoque estratégico}

La demostración es notablemente sencilla. Dado $\varepsilon>0$, usamos la
convergencia de $a$ para encontrar $N$ tal que $|a(m)-L|<\varepsilon$ para
todo $m\geq N$.

Para cualquier $n\geq N$, necesitamos $|(a\circ\sigma)(n)-L|<\varepsilon$.
La observación clave es que $\sigma(n)\geq n\geq N$ (usando el lema
$n\leq\sigma(n)$), de modo que $\sigma(n)\geq N$ y por tanto
$|a(\sigma(n))-L|<\varepsilon$.

\subsection{Demostración en Lean}

\begin{minted}{lean}
-- ---------------------------------------------------------------
-- Una subsucesión se obtiene aplicando a la sucesión original una
-- función de extracción; es decir, una función φ : ℕ → ℕ
-- estrictamente creciente. Por ejemplo, la subsucesión
--    u₀, u₂, u₄, u₆, ...
-- se obtiene con la función de extracción φ definida por
-- φ(n) = 2n.
--
-- Las definiciones anteriores se formalizan en Lean 4 como:
--
--    -- φ es una función de extracción.
--    def extraccion (φ : ℕ → ℕ) :=
--      StrictMono φ
--
--    -- v es una subsucesión de u.
--    def subsucesion (v u : ℕ → ℝ) :=
--      ∃ φ, extraccion φ ∧ v = u ∘ φ
--
--    -- a es el límite de u.
--    def LimSuc (u : ℕ → ℝ) (a : ℝ) :=
--      ∀ ε > 0, ∃ k : ℕ, ∀ n ≥ k, |u n - a| < ε
--
-- Demostrar que toda subsucesión de una sucesión convergente
-- converge al mismo límite que la sucesión.
-- ---------------------------------------------------------------


-- Demostración en lenguaje natural
-- ================================

-- Usaremos el siguiente lema: Si φ es una función de extracción,
-- entonces
--    ∀ n, n ≤ φ(n)
--
-- Por ser v una subsucesión de u, existe una función de
-- extracción φ tal que
--    v = u ∘ φ                                                (3)
--
-- Tenemos que demostrar que para cada ε > 0, existe un k ∈ ℕ tal
-- que
--     ∀ n ≥ k, |v(n) - a| < ε                                 (1)
--
-- Puesto que a es el límite de u, , existe un k ∈ ℕ tal
-- que
--     ∀ n ≥ k, |u(n) - a| < ε                                 (2)
--
-- Veamos que k verifica (1). Para ello, sea n ≥ k. Entonces,
--    φ(n) ≥ k                                                 (4)
-- ya que
--    φ(n) ≥ n    [por el Lema]
--         ≥ k
-- Luego,
--    |v(n) - a| = |(u ∘ φ )(n) - a|    [por (3)]
--               = |u(φ(n)) - a|
--               < ε                    [por (2) y (4)]

-- Demostraciones con Lean4
-- ========================

import Mathlib.Data.Real.Basic
import Analisis_en_Lean.C2_El_calculo_de_pi

open Nat

variable {u v : ℕ → ℝ}
variable {a : ℝ}
variable {φ : ℕ → ℕ}

-- φ es una función de extracción.
def extraccion (φ : ℕ → ℕ):=
  StrictMono φ

-- v es una subsucesión de u.
def subsucesion (v u : ℕ → ℝ) :=
  ∃ φ, extraccion φ ∧ v = u ∘ φ

-- 1ª demostración
-- ===============

example
  (hv : subsucesion v u)
  (ha : LimSuc u a)
  : LimSuc v a :=
by
  obtain ⟨φ, hφ, hφ'⟩ := hv
  -- φ : ℕ → ℕ
  -- hφ : extraccion φ
  -- hφ' : v = u ∘ φ
  intros ε hε
  -- ε : ℝ
  -- hε : ε > 0
  -- ⊢ ∃ k, ∀ n ≥ k, |v n - a| < ε
  obtain ⟨k, hk⟩ := ha ε hε
  -- k : ℕ
  -- hk : ∀ n ≥ k, |u n - a| < ε
  use k
  -- ⊢ ∀ n ≥ k, |v n - a| < ε
  intros n hn
  -- n : ℕ
  -- hn : n ≥ k
  -- ⊢ |v n - a| < ε
  have h1 : φ n ≥ k :=
    calc φ n
         ≥ n := StrictMono.le_apply hφ
       _ ≥ k := hn
  calc |v n - a|
       = |(u ∘ φ ) n  - a| := by rw [hφ']
     _ = |u (φ n) - a|     := rfl
     _ < ε                 := hk (φ n) h1

-- 2ª demostración
-- ===============

example
  (hv : subsucesion v u)
  (ha : LimSuc u a)
  : LimSuc v a :=
by
  obtain ⟨φ, hφ, rfl⟩ := hv
  -- φ : ℕ → ℕ
  -- hφ : extraccion φ
  -- ⊢ LimSuc (u ∘ φ) a
  intros ε hε
  -- ε : ℝ
  -- hε : ε > 0
  -- ⊢ ∃ k, ∀ n ≥ k, |(u ∘ φ) n - a| < ε
  obtain ⟨k, hk⟩ := ha ε hε
  -- k : ℕ
  -- hk : ∀ n ≥ k, |u n - a| < ε
  use k
  -- ⊢ ∀ n ≥ k, |(u ∘ φ) n - a| < ε
  intros n hn
  -- n : ℕ
  -- hn : n ≥ k
  -- ⊢ |(u ∘ φ) n - a| < ε
  have h1 : φ n ≥ k :=
    calc φ n
         ≥ n := hφ.le_apply
       _ ≥ k := hn
  calc |(u ∘ φ) n  - a|
     _ = |u (φ n) - a|     := rfl
     _ < ε                 := hk (φ n) h1

-- 3ª demostración
-- ===============

example
  (hv : subsucesion v u)
  (ha : LimSuc u a)
  : LimSuc v a :=
by
  obtain ⟨φ, hφ, rfl⟩ := hv
  -- φ : ℕ → ℕ
  -- hφ : extraccion φ
  -- ⊢ LimSuc (u ∘ φ) a
  intros ε hε
  -- ε : ℝ
  -- hε : ε > 0
  -- ⊢ ∃ k, ∀ n ≥ k, |(u ∘ φ) n - a| < ε
  obtain ⟨k, hk⟩ := ha ε hε
  -- k : ℕ
  -- hk : ∀ n ≥ k, |u n - a| < ε
  use k
  -- ⊢ ∀ n ≥ k, |(u ∘ φ) n - a| < ε
  intros n hn
  -- n : ℕ
  -- hn : n ≥ k
  -- ⊢ |(u ∘ φ) n - a| < ε
  apply hk
  -- ⊢ φ n ≥ k
  apply le_trans hn
  -- ⊢ n ≤ φ n
  exact hφ.le_apply

-- 4ª demostración
-- ===============

example
  (hv : subsucesion v u)
  (ha : LimSuc u a)
  : LimSuc v a :=
by
  obtain ⟨φ, hφ, rfl⟩ := hv
  -- φ : ℕ → ℕ
  -- hφ : extraccion φ
  -- ⊢ LimSuc (u ∘ φ) a
  intros ε hε
  -- ε : ℝ
  -- hε : ε > 0
  -- ⊢ ∃ k, ∀ n ≥ k, |(u ∘ φ) n - a| < ε
  obtain ⟨k, hk⟩ := ha ε hε
  -- k : ℕ
  -- hk : ∀ n ≥ k, |u n - a| < ε
  exact ⟨k, fun n hn ↦ hk (φ n) (le_trans hn hφ.le_apply)⟩

-- Lemas usados
-- ============

variable (m n x y z : ℕ)
#check (Nat.zero_le n : 0 ≤ n)
#check (le_trans : x ≤ y → y ≤ z → x ≤ z)
\end{minted}

\subsection{Demostración en lenguaje natural}

  Sea $\varepsilon>0$ dado. Como $a(n)\to L$, existe $N$ tal que para todo
  $m\geq N$:
  \[
    |a(m)-L| < \varepsilon
  \]

  Afirmamos que este mismo $N$ sirve para la subsucesión $a\circ\sigma$. Es
  decir, para todo $n\geq N$:
  \[
    |(a\circ\sigma)(n)-L| < \varepsilon
  \]

  Sea $n\geq N$. Por el lema \texttt{SubseqGe}, $n\leq\sigma(n)$, luego:
  \[
    \sigma(n)\geq n\geq N
  \]

  Como $\sigma(n)\geq N$:
  \[
    |(a\circ\sigma)(n)-L| = |a(\sigma(n))-L| < \varepsilon
  \]
  \hfill$\square$

\subsection{Por qué importa}

El teorema de la subsucesión tiene varias consecuencias importantes:
\begin{itemize}
  \item Si una sucesión tiene dos subsucesiones que convergen a límites
        distintos, la sucesión no converge. Esto proporciona una herramienta
        poderosa para probar divergencia.
  \item En el teorema de Bolzano-Weierstrass, toda sucesión acotada tiene
        una subsucesión convergente.
  \item Las subsucesiones se usan para extraer comportamiento «interesante»,
        como $\liminf$ y $\limsup$.
  \item El concepto se generaliza a espacios métricos y topológicos, donde
        juega un papel crucial en los argumentos de compacidad.
\end{itemize}

\section{Ejemplo de subsucesión: sucesiones oscilantes}

Aplicamos el teorema de la subsucesión a un ejemplo concreto: la sucesión
$a(n)=(-1)^n$, que oscila entre $-1$ y $1$. Esta sucesión no converge, pero
tiene subsucesiones convergentes.

Extrayendo los términos de índice par obtenemos una subsucesión constante que
converge a $1$. Análogamente, los términos de índice impar convergen a $-1$.

\subsection{Nuevas herramientas}

\paragraph{La táctica \texttt{let}.}
La táctica \texttt{let} es similar a \texttt{have}, pero crea un nuevo término
o función con nombre en lugar de solo demostrar una proposición. Para crear una
función en Lean se usa la palabra clave \texttt{fun}:
\begin{minted}{lean}
fun x ↦ expresion
\end{minted}
lo que define una función que toma la entrada $x$ y devuelve \texttt{expresión}.
Por ejemplo, \texttt{fun x ↦ x\^{}2} es la función cuadrática.

En el ejemplo de la subsucesión escribimos:

\begin{minted}{lean}
let σ : N → N := fun n ↦ 2 * n
\end{minted}
que define $\sigma(n)=2n$, seleccionando los términos de índice par.

\subsection{El enunciado matemático}

\begin{theorem}
  Si $a:\mathbb{N}\to\mathbb{R}$ está definida por $a(n)=(-1)^n$, entonces existe
  una subsucesión $\sigma$ y un límite $L$ tales que $\sigma$ es estrictamente
  creciente y $a\circ\sigma$ converge a $L$.
\end{theorem}

\subsection{Enfoque estratégico}

Definimos $\sigma(n)=2n$, que selecciona los índices pares: $0,2,4,6,\ldots$

Primero verificamos que $\sigma$ es estrictamente creciente: si $i<j$,
entonces $2i<2j$.

Luego calculamos $(a\circ\sigma)(n)=a(2n)=(-1)^{2n}=1$ para todo $n$.

Finalmente, esta sucesión constante converge trivialmente a $1$: para cualquier
$\varepsilon>0$, $|1-1|=0<\varepsilon$.

\subsection{Demostración en Lean}

\begin{minted}{lean}
-- ---------------------------------------------------------------

-- Demostración en lenguaje natural
-- ================================

-- Por el Reto 22, basta demostrar que la sucesión a tiene dos
-- subsucesiones que convergen a límites distintos. Para ello,
-- consideramos las funciones φ₁ y φ₂ de ℕ en ℕ definidas por
--    φ₁(n) = 2n
--    φ₂(n) = 2n + 1
-- Ambas son funciones de extracción ya que, para todo i, j ∈ ℕ
-- con i < j, se tiene que
--    φ₁(i) = 2i < 2j = φ₁(j)
--    φ₂(i) = 2i + 1 < 2j + 1 = φ₂(j)
-- Definimos las subsucesiones de a
--    b₁ = a ∘ φ₁
--    b₂ = a ∘ φ₂
-- Ambas subsucesiones son constantes (b₁ vale 1 y b₂ vale -1), ya
-- que para todo n ∈ ℕ
--    b₁(n) = a(2n) = (-1)²ⁿ = 1
--    b₂(n) = a(2n + 1) = (-1)²ⁿ⁺¹ = -1
-- Luego, aplicando el Reto 23 que dice "toda sucesión constante
-- converge a su valor", b₁ converge a 1 y b₂ converge a -1. Como
-- 1 ≠ -1, a tiene dos subsucesiones con límites distintos y, por
-- el Reto 22, no es convergente.

-- Demostraciones en Lean 4
-- ========================

import Mathlib.Tactic
import RetosLean4.Defs -- LimSuc, SucConvergente, extraccion, subsucesion
import RetosLean4.Reto_9 -- unicidad_limite
import RetosLean4.Reto_21 -- limite_subsucesion
import RetosLean4.Reto_22 -- no_convergencia_subsucesiones
import RetosLean4.Reto_23 -- limSuc_const

variable {a : ℕ → ℝ}

-- 1ª demostración
-- ===============

lemma extraccion_par : extraccion (fun n => 2 * n) :=
by
  intro i j hij
  -- i j : ℕ
  -- hij : i < j
  -- ⊢ (fun n => 2 * n) i < (fun n => 2 * n) j
  dsimp
  -- ⊢ 2 * i < 2 * j
  omega

lemma extraccion_impar : extraccion (fun n => 2 * n + 1) :=
by
  intro i j hij
  -- i j : ℕ
  -- hij : i < j
  -- ⊢ (fun n => 2 * n + 1) i < (fun n => 2 * n + 1) j
  dsimp
  -- ⊢ 2 * i + 1 < 2 * j + 1
  omega

lemma suc_par
  (ha : ∀ n, a n = (-1)^n)
  (n : ℕ)
  : a (2 * n) = 1 :=
by
  calc a (2 * n)
         = (-1) ^ (2 * n) := ha (2 * n)
       _ = 1              := by simp

lemma suc_impar
  (ha : ∀ n, a n = (-1)^n)
  (n : ℕ)
  : a (2 * n + 1) = -1 :=
by
  calc a (2 * n + 1)
         = (-1) ^ (2 * n + 1) := ha (2 * n + 1)
       _ = -1                 := by simp [pow_succ]

example
  (ha : ∀ n, a n = (-1)^n)
  : ¬SucConvergente a :=
by
  let φ₁ := fun n : ℕ => 2 * n
  let φ₂ := fun n : ℕ => 2 * n + 1
  have hb₁ : subsucesion (a ∘ φ₁) a := ⟨φ₁, extraccion_par, rfl⟩
  have hb₂ : subsucesion (a ∘ φ₂) a  := ⟨φ₂, extraccion_impar, rfl⟩
  have hL₁ : LimSuc (a ∘ φ₁) 1 :=
    limSuc_const (fun n => suc_par ha n)
  have hL₂ : LimSuc (a ∘ φ₂) (-1) :=
    limSuc_const (fun n => suc_impar ha n)
  have h : (1:ℝ) ≠ -1 := by norm_num
  exact no_convergencia_subsucesiones hb₁ hL₁ hb₂ hL₂ h

-- 2ª demostración
-- ===============

example
  (ha : ∀ n, a n = (-1)^n)
  : ¬SucConvergente a :=
by
  have h1 : extraccion (fun n => 2 * n) := extraccion_par
  have hb₁ : subsucesion (a ∘ fun n => 2 * n) a := ⟨_, h1, rfl⟩
  have hL₁ : LimSuc (a ∘ fun n => 2 * n) 1 :=
    limSuc_const (fun n => suc_par ha n)
  have h2 : extraccion (fun n => 2 * n + 1) := extraccion_impar
  have hb₂ : subsucesion (a ∘ fun n => 2 * n + 1) a := ⟨_, h2, rfl⟩
  have hL₂ : LimSuc (a ∘ fun n => 2 * n + 1) (-1) :=
    limSuc_const (fun n => suc_impar ha n)
  exact no_convergencia_subsucesiones hb₁ hL₁ hb₂ hL₂ (by norm_num)

-- 3ª demostración
-- ===============

theorem no_convergente_pot_menos_uno
  (ha : ∀ n, a n = (-1)^n)
  : ¬SucConvergente a :=
no_convergencia_subsucesiones
  ⟨_, extraccion_par, rfl⟩
  (limSuc_const (fun n => suc_par ha n))
  ⟨_, extraccion_impar, rfl⟩
  (limSuc_const (fun n => suc_impar ha n))
  (by norm_num)
\end{minted}

\subsection{Demostración en lenguaje natural}

  \textbf{Paso 1: $\sigma$ es una subsucesión.}

  Sean $i,j\in\mathbb{N}$ con $i<j$. Entonces:
  \[
    \sigma(i) = 2i < 2j = \sigma(j)
  \]
  Por tanto $\sigma$ es estrictamente creciente, y en consecuencia es una
  subsucesión.

  \medskip
  \textbf{Paso 2: $a\circ\sigma$ converge a $1$.}

  Sea $\varepsilon>0$ dado. Para todo $n\geq 0$:
    \begin{align*}
      (a\circ\sigma)(n) &= a(\sigma(n)) \\
               &= a(2n)   \\
               &= (-1)^{2n} \\
               &= \bigl((-1)^2\bigr)^n \\
               &= 1^n \\
               &= 1
    \end{align*}
  Por tanto:
    \begin{align*}
      |(a\circ\sigma)(n)-1| &= |1-1| \\
                   &= 0 \\
                   &< \varepsilon
    \end{align*}
  \hfill$\square$

\subsection{Implicaciones para la divergencia}

Este ejemplo demuestra que $a(n)=(-1)^n$ no converge. Podemos extraer
subsucesiones que convergen a $1$ (índices pares) y a $-1$ (índices impares).
Como una sucesión convergente debe tener todas sus subsucesiones convergiendo
al mismo límite, y hemos encontrado dos con límites distintos, la sucesión $a$
no puede converger.

Esta técnica —mostrar la divergencia exhibiendo dos subsucesiones con límites
distintos— es una herramienta poderosa en el análisis.

\section{Ejercicios}



%%% Local Variables:
%%% mode: LaTeX
%%% TeX-master: "Analisis_en_Lean"
%%% End:
